\documentclass[11pt]{article}
\usepackage[margin=1in]{geometry}
\usepackage{amsmath,amssymb,amsthm}
\usepackage{booktabs}
\usepackage{hyperref}
\emergencystretch=4em
\hypersetup{colorlinks=true,linkcolor=blue,urlcolor=blue}

\newcommand{\eps}{\varepsilon}
\newcommand{\epsth}{\eps_\theta}
\newcommand{\epsfr}{\eps_{\bar\theta}}
\newcommand{\xs}{x_\sigma}
\newcommand{\E}{\mathbb{E}}
\newcommand{\R}{\mathbb{R}}
\newcommand{\norm}[1]{\left\lVert #1 \right\rVert}
\newcommand{\cnull}{\varnothing}
\newcommand{\pmi}{\mathrm{pmi}}

\theoremstyle{plain}
\newtheorem{proposition}{Proposition}
\newtheorem{corollary}[proposition]{Corollary}
\theoremstyle{remark}
\newtheorem{remark}[proposition]{Remark}

\title{The guided objective the repository trains and the unguided one the board writes,\\
reconciled in one notation}
\author{poe\_repair\_min, thread off \texttt{designing-the-correction-loss}}
\date{2026-09-16}

\begin{document}
\maketitle

\begin{abstract}
The trainer fits a guided composition to a guided joint-prompt target at $w = 7.5$; the
supervisor's board writes the same fit with no guidance anywhere. The two agree exactly when the
adapted null branch equals the frozen one, and diverge in proportion to the null branch's drift,
amplified by $(w-1)$. Measured on the four thirty-thousand-step arms, the amplified drift term is
$67$ times the entire guided residual in the unanchored run and $0.95$ times it in the anchored
one, so the plain adapter provably does \emph{not} solve the unguided problem while the anchored
one is merely not ruled out. The document also settles that no assignment of three branch
functions can match the joint target without moving the single-concept branches, names the three
candidate protection terms and which one to add, and gives one cached-data probe for the
off-policy gap.
\end{abstract}

\tableofcontents

\section{Setup, and one notation for both objectives}

\paragraph{Companion document.}
This document does not re-derive the additive skeleton, the gauge proposition, or the step-0
floor. Those are in \texttt{correction-loss-variants.tex} beside this file and are used here by
reference.

\paragraph{States, branches, frozen and adapted.}
Write $\xs$ for a noised latent at scale $\sigma$, $\epsth(\xs,\sigma,c)$ for the adapted
noise prediction and $\epsfr$ for the frozen one. Fix a state and a pair $(c_1,c_2)$ and set
\begin{equation}
  u \;:=\; \epsth(\xs,\sigma,\cnull), \qquad
  \bar u \;:=\; \epsfr(\xs,\sigma,\cnull), \qquad
  \delta \;:=\; u - \bar u ,
  \label{eq:coords}
\end{equation}
with $a := \epsth(\xs,\sigma,c_1)$, $b := \epsth(\xs,\sigma,c_2)$, and $\eps_J :=
\epsfr(\xs,\sigma,c_{1\wedge 2})$ the frozen joint-prompt prediction read from cache. The
unguided composition and its guided form are
\begin{equation}
  T \;=\; a + b - u , \qquad
  T^{(w)} \;=\; u + w(a-u) + w(b-u) \;=\; w\,T - (w-1)\,u ,
  \label{eq:compose}
\end{equation}
implemented as \texttt{poe\_eps} and \texttt{guided\_eps}
(\texttt{poe\_repair/\_sdxl/metrics.py:12--17}) and combined at
\texttt{trainer.py:331--332} with $w = 7.5$. The target is guided against the \emph{frozen}
null, $J^{(w)} := \bar u + w(\eps_J - \bar u) = w\,\eps_J - (w-1)\,\bar u$
(\texttt{trainer.py:388}).

\paragraph{The two objectives.}
\begin{align}
  \mathcal{L}_A &\;=\; \norm{\,T^{(w)} - J^{(w)}\,}^2
  \qquad\text{(the repository, \texttt{trainer.py:450--453})},
  \label{eq:LA}\\
  \mathcal{L}_{B2} &\;=\; \norm{\,T - \eps_J\,}^2
  \qquad\text{(the board, mono-as-target variant)},
  \label{eq:LB2}\\
  \mathcal{L}_{B1} &\;=\; \norm{\,z - \big(a + b - u\big)\,}^2,
  \quad x \sim p_{\mathrm{co}},\ \xs = x + \sigma z
  \qquad\text{(the board, data-as-target variant)}.
  \label{eq:LB1}
\end{align}
Write $\Delta := T - \eps_J$ for the unguided residual, so $\mathcal{L}_{B2} = \norm{\Delta}^2$.
All three carry a per-level weight $\lambda(\sigma)$; the runs discussed here use
$\lambda \equiv 1$ (\texttt{loss\_eps} equals \texttt{loss\_fit} to five figures on every arm),
so every number below is in one unit, the per-element mean square of a noise-space tensor.

\paragraph{The identity this document turns on.}
From \eqref{eq:compose},
\begin{equation}
  \boxed{\;T^{(w)} - J^{(w)} \;=\; w\,\Delta \;-\; (w-1)\,\delta\;}
  \label{eq:key}
\end{equation}
so
\begin{equation}
  \mathcal{L}_A \;=\; w^2\,\mathcal{L}_{B2}
  \;-\; 2w(w-1)\,\langle \Delta, \delta\rangle
  \;+\; (w-1)^2 \norm{\delta}^2 .
  \label{eq:expand}
\end{equation}
At initialisation the adapter is zero, $\delta = 0$, and $\mathcal{L}_A = w^2 \mathcal{L}_{B2}$
exactly; the numerical check on one cached cell returned $56.2500$ against $w^2 = 56.25$.
Everything that follows is about the two terms that switch on when $\delta$ moves.

\paragraph{The runs these numbers come from.}
Four arms under
\texttt{/datasets/mmolefe/poe\_repair\_min/}\allowbreak\texttt{outputs/correction\_loss\_variants/}, each rank $16$,
$30{,}000$ optimizer steps, denoising steps $0$ to $25$, $1024^2$, $w = 7.5$, pair
\texttt{a cat} $\times$ \texttt{a dog}. They differ in \texttt{branch\_prompt\_style}
(\texttt{plain}, \texttt{plurality}, \texttt{connective}) and, for the fourth,
\texttt{null\_anchor} $ = 10.0$ against $0.0$ elsewhere (verified in each arm's
\texttt{config.json}).

\section{Objective 1: is matching the unguided composition to the mono prediction well posed?}

\subsection{The density reading is false, and it is false exactly where the project works}

\begin{proposition}[What exact matching would demand]
\label{prop:pmi}
Suppose the three branches were the true scores of $p(x\mid c_1)$, $p(x\mid c_2)$ and $p(x)$, and
the target were the true score of $p(x \mid c_1, c_2)$. Then
\begin{equation}
  \Delta \;=\; \sigma \nabla_x \Big[\log \tfrac{p(c_1,c_2\mid x)}{p(c_1\mid x)p(c_2\mid x)}\Big]
  \;=\; \sigma \nabla_x\, \pmi_x(c_1;c_2),
\end{equation}
so $\Delta \equiv 0$ if and only if the pointwise mutual information between the two concepts
given the image is constant in $x$. Conditional independence given $x$ is sufficient but not
necessary; a constant offset is invisible to the score.
\end{proposition}
\begin{proof}
Bayes on each of the three conditionals turns the ratio
$p(x\mid c_1,c_2)\big/\!\left[p(x\mid c_1)p(x\mid c_2)/p(x)\right]$ into
$\exp(\pmi_x - \pmi)$, with $\pmi$ the prior-level term, free of $x$. Take $\log$ and
$\nabla_x$; the prior term drops. The sign convention matches this project's
\[
  r_t \;=\; -\sigma_t \nabla_{x_t}\log\!\left[\frac{p(c_1,c_2\mid x_t)}{p(c_1\mid x_t)\,p(c_2\mid x_t)}\right]
\]
in \texttt{artifacts/notes/interaction-term-as-pmi-gradient/}.
\end{proof}

\paragraph{Why that condition cannot hold here.}
$\pmi_x$ is near zero on an image showing a cat and a dog as separate objects (both posteriors
near one, the joint posterior near one) and strongly negative on a chimera (both single-concept
posteriors moderate, the conjunction posterior low). The gradient of $\pmi_x$ is therefore
nonzero precisely on the region of state space where composition fails, which is the region the
adapter is trained on. \emph{Confidence: the derivation is confident; the claim about the two
posteriors' behaviour is a statement about the scorer's own semantics and is likely but
unverified, since no run has measured $p(c_1,c_2\mid x)$ against $p(c_1\mid x)p(c_2\mid x)$ on
this pool.}

\subsection{So what is actually being asked of the three branches}

The regression does not ask the model to make two concepts independent, and it cannot: the network
does not touch $p(x \mid c_1, c_2)$. What it asks is narrower and should be stated that way in the
paper:

\begin{quote}
Find three functions $a(\cdot,c_1)$, $b(\cdot,c_2)$, $u(\cdot)$, each reading one prompt, whose
fixed linear combination $a+b-u$ reproduces a fourth function, $\epsfr(\cdot,c_{1\wedge2})$, on the
state distribution of plain product-of-experts sampling trajectories.
\end{quote}

Three separate things are being conflated when this is described as removing an independence
assumption. The target is the model's joint-\emph{prompt} branch, not the true joint density; the
joint prompt string is not the conjunction event (this is unstated premise A of the pressure-test
note, and it is the first thing a reviewer will find); and the fit is over one state distribution,
not over $\R^d$.

\subsection{A solution exists, and it cannot leave the single-concept branches alone}

\begin{proposition}[No damage-free solution]
\label{prop:damage}
If every branch keeps its frozen value, $a = \epsfr(\cdot,c_1)$, $b = \epsfr(\cdot,c_2)$,
$u = \bar u$, then $\Delta = \epsfr(c_1)+\epsfr(c_2)-\bar u - \eps_J$, which is the cached
correction the project calls $r_t$ and is nonzero by construction. Hence any $\theta$ with
$\Delta = 0$ moves at least one branch off its frozen value.
\end{proposition}

This is a tautology, and it is worth writing down because it converts objective 1 from a yes/no
question into a magnitude question: not \emph{can} single-concept behaviour survive, but \emph{how
far must each branch move}. Writing $\phi_i$ for the per-concept displacement and $m$ for the shared
one, the additive skeleton of the companion document requires $m + \phi_i + \phi_j = \Delta_{ij}$
at each state, and the step-0 fit already on disk says such a $\phi$ exists to within
\begin{center}
\small
\begin{tabular}{lrrrr}
\toprule
 & in-sample & held out & shuffled ceiling & algebraic control \\
\midrule
seed 1, 1003 pairs, 294 concepts & $0.0280$ & $0.0884$ & $0.2166$ & $0.00022$ \\
seed 2, 1010 pairs, 299 concepts & $0.0249$ & $0.0772$ & $0.2329$ & $0.00014$ \\
\bottomrule
\end{tabular}
\end{center}
as a fraction of the correction's squared energy. The in-sample and held-out columns are
the five-fold split of \texttt{additive\_transfer\_step0.json}, fitted and evaluated on the same
folds. The shuffled ceiling and the algebraic control are from the all-pairs fit in
\texttt{additive\_floor\_step0.json}, whose own in-sample floor is $0.0315$ and $0.0278$. Both
files sit under \texttt{outputs/idea\_probes/}. The held-out column is the one that licenses the approach: per-concept
functions fitted on $80\%$ of pairs predict the remaining $20\%$ to within $8\%$ of the correction,
so the residual is genuinely per-concept and not pair memorisation.

\paragraph{The number this section needs and does not have.}
The size of the damage, $\norm{\phi_i}$ relative to the branch's own guidance signal
$\norm{\epsfr(\cdot,c_i) - \bar u}$, at step 0. The fit that produces $\phi$ is already computed
inside \texttt{scripts/idea\_probes/additive\_floor\_step0.py}; it solves for the per-concept
factors and then discards them, reporting only the residual. Recording
$\norm{\phi_i}/\norm{\epsfr(c_i)-\bar u}$ per concept is one added line and no new forward passes.
Until that number exists, ``the correction is small'' is a claim about the residual, not about the
distortion. \emph{Confidence: family-only on the magnitude; confident that the quantity is
computable from what is on disk.}

\paragraph{Verdict on objective 1.}
Well posed as a regression, not well posed as a density statement. A solution exists and is
reachable by per-concept functions (held-out $8\%$), and it necessarily moves the single-concept
branches; the independence reading should not appear in the paper. The saving grace is
distributional rather than algebraic: the branches are only constrained on pair-trajectory states,
so single-concept behaviour is untouched \emph{off} those states, neither damaged nor protected.
That is the same fact as objective 4 read from the other side.

\section{Objective 2: does the null anchor make the guided and unguided objectives agree?}

\subsection{The exact statement, which is about freezing and not about anchoring}

\begin{proposition}[Hard freeze gives identity at every $w$]
\label{prop:freeze}
If the null branch is held at its frozen value as a \emph{function} of $\theta$, that is
$u(\theta) \equiv \bar u$ so $\partial u/\partial\theta = 0$, then by \eqref{eq:key}
$\mathcal{L}_A = w^2\,\mathcal{L}_{B2}$ and $\nabla_\theta \mathcal{L}_A =
w^2\,\nabla_\theta \mathcal{L}_{B2}$, identically in $\theta$ and for every $w$. The guidance
weight becomes a constant loss scale, absorbable into the learning rate, and one trained adapter is
simultaneously optimal unguided and at $w = 7.5$.
\end{proposition}

\begin{corollary}[The freeze is free in capacity]
By the gauge proposition of the companion document, a free null branch and a frozen one reach
exactly the same set of composed predictions $T$. So freezing costs nothing in the unguided
problem, removes the sampler-amplified unidentified direction, and drops the null forward from the
training step (the frozen $\eps_\cnull$ is already cached as \texttt{eps\_uncond}), cutting the
$3K$-wide forward at \texttt{trainer.py:440--450} to $2K$.
\end{corollary}

\begin{remark}[This is a claim about a hard freeze, and the repository implements a soft penalty]
\texttt{trainer.py:496--503} adds $\mu\norm{\delta}^2$ with $\mu = 10$ on the anchored arm. That
leaves $\delta \neq 0$, so Proposition~\ref{prop:freeze} does not apply and the cross term and the
quadratic term of \eqref{eq:expand} both survive. Whether they matter is an empirical question with
an answer already logged.
\end{remark}

\subsection{Where it breaks, measured}

Both \texttt{loss\_fit} ($=\mathcal{L}_A$) and \texttt{loss\_null\_anchor} ($=\norm{\delta}^2$) are
logged in the same units, per-element mean square, at the same step (\texttt{trainer.py:505--514}).
The amplified drift term of \eqref{eq:expand} is therefore $(w-1)^2 \times
\texttt{loss\_null\_anchor} = 42.25 \times \texttt{loss\_null\_anchor}$, directly comparable to
\texttt{loss\_fit}. At step $30{,}000$:

\begin{center}
\begin{tabular}{lrrrr}
\toprule
arm & $\mu$ & $\mathcal{L}_A$ & $(w-1)^2\norm{\delta}^2$ & ratio \\
\midrule
\texttt{plain}      & $0$    & $4.9\times10^{-4}$ & $3.30\times10^{-2}$ & $67.3$ \\
\texttt{plurality}  & $0$    & $6.0\times10^{-4}$ & $2.32\times10^{-2}$ & $38.7$ \\
\texttt{connective} & $0$    & $5.0\times10^{-4}$ & $3.34\times10^{-2}$ & $66.8$ \\
\texttt{anchor}     & $10$   & $8.9\times10^{-4}$ & $8.45\times10^{-4}$ & $\mathbf{0.95}$ \\
\bottomrule
\end{tabular}
\end{center}

\paragraph{Reading the unanchored rows.}
A ratio of $67$ says the null drift's contribution to the guided residual is sixty-seven times the
whole guided residual. The two terms of \eqref{eq:key} are therefore cancelling: $w\Delta$ is large
and $(w-1)\delta$ is large and they nearly annihilate. That is the gauge degeneracy behaving exactly
as predicted, not a coincidence. By $\norm{w\Delta} \ge \norm{(w-1)\delta} - \norm{T^{(w)}-J^{(w)}}$,
\begin{equation}
  \mathcal{L}_{B2}^{\ \texttt{plain}} \in [4.52\times10^{-4},\ 7.38\times10^{-4}],
  \qquad
  \mathcal{L}_A/w^2 = 8.71\times10^{-6},
  \label{eq:plainbracket}
\end{equation}
so the plain arm's unguided loss is between $52$ and $85$ times $\mathcal{L}_A/w^2$. \textbf{The
plain adapter is a certificate of failure on the unguided problem}: whatever it learned, it did not
learn a composition that works at $w=1$. \emph{Confidence: confident. The bracket is a triangle
inequality on two logged quantities in matching units; nothing is assumed.}

\paragraph{Reading the anchored row.}
The same bracket for the anchored arm is $\mathcal{L}_{B2} \in [1.0\times10^{-8},\
6.2\times10^{-5}]$ against $\mathcal{L}_A/w^2 = 1.58\times10^{-5}$, a ratio bracket of
$[0.0007,\ 3.9]$. It contains $1$. So the anchor takes the divergence from certified to
undetermined, and no more than that. At $\mu = 10$ the drift's amplified contribution is still the
same size as the entire residual it perturbs, which is the honest version of the claim.

\paragraph{The number that would settle it.}
$\norm{\Delta}^2$ logged directly, as a per-element mean square, beside \texttt{loss\_fit}. Both
$T$ (unguided) and $\eps_J$ (cached raw) are in hand at \texttt{trainer.py:450--453} before the
guided forms are built, so this is one line under \texttt{no\_grad} and no extra forward pass. With
$\mathcal{L}_{B2}$ logged, \eqref{eq:expand} gives the cross term by subtraction and the question
stops being a bracket.

\paragraph{Verdict on objective 2.}
The claim is right in direction and, as implemented, not established. Driving $\delta$ to zero does
make the two objectives agree at every $w$, and the correct way to do it is to freeze the null
branch rather than penalise its drift: the freeze is exact, costs no capacity, and saves a third of
the training forward. The soft anchor at $\mu = 10$ reduces the amplified drift term by a factor of
$39$ but leaves it at parity with the residual, so a single anchored adapter being correct at both
$w = 1$ and $w = 7.5$ is currently unfalsified rather than shown.

\begin{remark}[The freeze must be matched at sampling time]
The adapter is LoRA on cross-attention projections, so it modifies the null branch too: the empty
prompt still produces keys and values. Training with a frozen null while sampling with an adapted
null is a train-test mismatch that reintroduces exactly the direction the freeze removed. The
sampler must run its $\cnull$ forward with the adapter detached. This repository has already paid
for the general version of this mistake once, in a windowed sampler that left the adapter enabled
while rendering its own references.
\end{remark}

\section{Objective 3: the three candidate single-concept protections}

\paragraph{(i) The null anchor, implemented.}
$\mu \norm{\epsth(\cnull) - \epsfr(\cnull)}^2$ on pair-trajectory states, \texttt{trainer.py:496}.
Free, because $\epsfr(\cnull)$ is cached. \textbf{It is not single-concept protection and should
not be described as such.} Its job is gauge-fixing, per Section 3; it constrains one branch, the one
that carries no concept.

\paragraph{(ii) The board's data term.}
$\E_{x \sim p_{\mathrm{single}},\,z}\norm{z/\sigma - s_\theta(x+\sigma z,\sigma,c_{\mathrm{dog}})}^2$.
This anchors the dog branch to the \emph{true} data score at states drawn from noised real dog
images. Two consequences follow from that word ``true''. It is a fine-tuning term, not a
preservation term: where the frozen model was already wrong about dogs, this moves the branch to a
third place, neither frozen nor current. And it needs real single-concept images, which is a much
weaker data requirement than $\mathcal{L}_{B1}$'s (one concept per image, not two in one image):
the Open Images coverage probe matched $83$ of $317$ pool concepts exactly, so this term is
reachable for roughly a quarter of the vocabulary and unreachable for the rest.

\paragraph{(iii) The board's frozen-versus-trained term.}
$\norm{\epsfr(\cdot,c_{\mathrm{dog}}) - \epsth(\cdot,c_{\mathrm{dog}})}^2$. This is function-space
distillation of the model onto itself, and it is the null anchor generalised from $c = \cnull$ to
$c = c_i$. It needs no images at all, only states, and the states are a free choice. That freedom is
the whole design question, because:

\begin{proposition}[Applying (iii) on the training states cancels the fit]
\label{prop:conflict}
On pair-trajectory states, the fit requires $\phi_i \neq 0$ (Proposition~\ref{prop:damage}) and
(iii) penalises $\norm{\phi_i}^2$. In the limit of a large weight, all three branches are pinned to
frozen, $T = \bar T$, and the correction is identically zero. So (iii) evaluated where the fit lives
does not protect single-concept behaviour; it prices the correction.
\end{proposition}

\paragraph{Recommendation: (iii), on single-concept trajectory states.}
Build a small cache of states from single-prompt sampling runs (one branch per step, no pair, no
training), and evaluate (iii) there. This separates the two jobs cleanly: correct the composition
where pairs go, preserve the model where a single concept goes. The costs are one sampling run per
protected concept, once, and one extra adapted forward per training step per protected concept; the
frozen prediction at those states is cached alongside, exactly as \texttt{eps\_uncond} is today, so
the penalty side is free at train time.

Ranking, with reasons rather than a tie-break: (iii) on single-concept states is the only one of the
three that both targets the concept branches and does not fight the fit. (i) is worth keeping
regardless, because it does a different job (gauge), and Section 3 recommends replacing it with the
stronger hard freeze. (ii) is a different project: it improves the model against data rather than
preserving it, and it is available for a quarter of the vocabulary.

\paragraph{What would show (iii) is doing something, and what would show it is too strong.}
Sweep its weight. The correction's size \texttt{delta\_hat\_norm} against the target's
\texttt{delta\_target\_norm} ($45.12$ on these cells) reads the price; the pair compose rate on the
held-out seeds reads whether the price was worth paying; a single-concept render set scored by the
same instance-count scorer reads what was bought. \emph{Confidence: likely but verify. The conflict
of Proposition~\ref{prop:conflict} is a derivation; that single-concept trajectory states are
sufficiently disjoint from pair-trajectory states to make the separation real is a conjecture, and
it is testable for free by measuring the distance between the two state clouds at each step.}

\section{Objective 4: what the objective does not guarantee off the cached trajectories}

\paragraph{What it does guarantee.}
Every cached $\xs$ came from a plain product-of-experts sampling run with no adapter. So the fit is a
statement in $L^2(\rho_{\mathrm{PoE}})$: the composed prediction is close to the target in mean
square under the state distribution that plain PoE visits, at denoising steps $0$ to $25$ and at the
seeds and pairs in the cache. Nothing more. In particular it says nothing at states the adapter
itself reaches, and the adapter changes the trajectory from the first step it acts on, so those are
never the same states after step $0$.

\paragraph{What that costs.}
This is behaviour cloning under distribution shift, and it fails the standard way: per-step error
compounds along the rollout, because each step's state error feeds the next step's prediction error.
A bound on the fit at cached states does not bound the deviation at the end of a fifty-step
trajectory. \emph{Confidence: confident on the structure, family-only on the magnitude here, since
no measurement in this repository isolates the compounding term.}

\paragraph{One piece of evidence already on disk, and what it is not.}
\texttt{scripts/}\allowbreak\texttt{trajectory\_divergence.py} measures, per step, the CLIP and DINOv2 cosine distance
between decoded frames of the correction-on and correction-off trajectories. The curves separate
early: steepest rise at step $5$ or $6$ in three of the cat-dog cells, later (steps $22$ to $23$)
in others. Its output is \texttt{divergence.json}, written under the interaction-term
outputs in \texttt{cache\_analyses/}\allowbreak\texttt{trajectory\_divergence/}. This is
the oracle correction, not the trained adapter, so it bounds the question from a neighbouring
direction rather than answering it: it says a correction of this kind moves the trajectory to a
visibly different place well inside the trained window.

\subsection{The probe: one-step-off-trajectory loss}

Uses only the existing cache, one trained checkpoint, no training and no rendering.

\begin{enumerate}
  \item For each cached step $k \in \{0,\dots,25\}$ and each cached cell, run the trained adapter's
    three branches at the cached $x_k$, compose at $w = 7.5$, and take one DDIM step to get
    $x'_{k+1}$. The cached $x_{k+1}$ is the plain-PoE successor of the same state.
  \item \textbf{Coverage, at zero forward cost.} Report
    $d_k = \norm{x'_{k+1} - x_{k+1}}$ against $s_k$, the standard deviation of the cached $x_{k+1}$
    across the eight seeds of that pair. $d_k \ll s_k$ means the cached cloud already covers where
    the adapter goes in one step.
  \item \textbf{The loss gap.} Evaluate $\mathcal{L}_A$ at $x'_{k+1}$ and at $x_{k+1}$. The
    displaced state needs a fresh frozen joint-prompt forward, since the cached $\eps_J$ is keyed to
    the cached state: one frozen forward per probe state, a few hundred in total, minutes on one
    card.
  \item \textbf{The bar, written here before the run.} If
    $\mathcal{L}_A(x'_{k+1}) \le 2\,\mathcal{L}_A(x_{k+1})$ at every $k$ in the trained window, the
    fit survives one step of its own drift and off-policy shift is not the binding problem. If the
    ratio exceeds $2$ anywhere in the window, the training distribution must be refreshed against the
    adapter (re-cache with the adapter attached, iterated), and no amount of further training on the
    frozen cache will fix it.
\end{enumerate}

Only step 3 costs a forward pass; steps 1 and 2 are tensor arithmetic on cached data. The bar belongs
in the script's source before the number exists, matching this repository's convention for
\texttt{MIN\_MEDIAN\_RATIO}.

\section{Objective 5: is \texorpdfstring{$\mathcal{L}_{B1}$}{the data-target objective} reachable here?}

\paragraph{Verdict: no, and not for lack of effort in the search.}
$\mathcal{L}_{B1}$ needs $p_{\mathrm{co}}$, real images with both concepts present as separate
objects. This repository holds no real image data: \texttt{data/pilot/} is SDXL output and no dataset
loader exists in \texttt{poe\_repair/} or \texttt{scripts/}. Against Open Images' $601$ boxable
classes, $83$ of $317$ pool concepts match exactly and $11$ of $1146$ pairs have at least one
co-occurring validation image (\texttt{outputs/idea\_probes/openimages\_coverage.json}).

The eleven are the finding, not the shortage. They are whale with dolphin, swan with duck, owl with
falcon, turtle with tortoise: pairs that share a habitat, look alike, and already sit in the easy
tail of the additive fit. The mechanism is not an accident of one corpus. A pair has co-occurrence
photographs because the two things genuinely turn up together, and that same co-occurrence in the
training corpus is why the model already composes them. The data is abundant exactly where it is not
needed and absent exactly where it is. \emph{Confidence: confident for Open Images, which is
measured; likely but verify as a general statement about photographic corpora, since only one corpus
has been checked.}

\paragraph{The nearest reachable objective, first choice: $\mathcal{L}_{B2}$, which is already running.}
Up to the $w^2$ factor and the null drift of Section 3, the board's mono-target variant \emph{is} what
the repository trains. Freezing the null branch closes the gap exactly (Proposition~\ref{prop:freeze}),
at which point the board's second variant and the repository's objective are the same objective and
the disagreement dissolves. Nothing new is needed.

\paragraph{Second choice, and the one worth designing: a rejection-sampled corpus from the model itself.}
Replace $p_{\mathrm{co}}$ with the model's own samples that pass a composition filter. This is
reachable today: searching eight nearby initial noises composes $8$ of $8$ on the cat-dog held-out
seeds where the adapter alone composes $6$ of $8$, by the validated instance count
(\texttt{report/is-the-gap-the-samplers-or-the-models/}\allowbreak
\texttt{does-searching-over-the-noise-sharpen-the-corrected-render.md}). Images in which both
animals appear as separate objects can therefore be manufactured for pairs that fail. The corpus becomes noised real-ish images from
the model's own successful tail, and $\mathcal{L}_{B1}$ runs on it unchanged.

Two things make this more than self-distillation with extra steps. The current target is known to be
unreliable: scope 08 drops training cells whose joint target shows the wrong animals, and names
``whether a better teacher than the joint prompt would beat all of this'' as its own follow-on
(\texttt{plans/08-improve-the-rank-32-pooled-adapter/}\allowbreak\texttt{review/04-the-parent-and-the-four-children.md},
lines 39 and 147). A filtered sample set is a teacher that is right by construction where the joint
prompt is wrong. Against that, the filter is the risk: this project's validated instance-count scorer
counts instances and cannot tell a cat beside a dog from two dogs, so the corpus needs the second
head that reads which two animals are present, or it trains the composition toward a density of
duplicated single concepts. \emph{Confidence: likely but verify. The compose counts are measured; that
a filtered corpus trains better than the joint-prompt target is a conjecture and is the experiment,
not the premise.}

\paragraph{Third, unsettled: caption retrieval over a web-scale corpus.}
Filtered by the same scorer. It needs no fixed class vocabulary, which is what eliminated Open Images.
This remains the open check of the idea map's claim 6 and nothing here settles it.

\section{The four verdicts}

\begin{enumerate}
  \item \textbf{Unguided-to-mono target.} Well posed as a regression, not as a density statement.
    Asking $a+b-u = \eps_J$ asks for $\nabla_x \pmi_x(c_1;c_2) = 0$, which is false exactly where
    composition fails, so the adapter is not learning independence; it is fitting three
    single-prompt functions to reproduce a fourth under a fixed linear rule on one state
    distribution. A solution exists (held-out per-concept residual $8\%$ of the correction at step
    0) and provably moves the single-concept branches. \emph{Confident.}
  \item \textbf{Null anchor and the two objectives.} A \emph{hard freeze} of the null branch makes
    $\mathcal{L}_A = w^2 \mathcal{L}_{B2}$ identically at every $w$, costs no capacity, and saves a
    third of the training forward. The implemented \emph{soft} anchor does not: at $\mu = 10$ the
    amplified drift is still $0.95$ times the entire residual, against $67$ times unanchored. The
    plain arm is certified not to solve the unguided problem ($52$ to $85$ times worse); the
    anchored arm is undetermined. \emph{Confident on the derivation and on both measured ratios.}
  \item \textbf{Single-concept protection.} Add the frozen-versus-trained term, evaluated on states
    from single-concept sampling trajectories rather than on the pair trajectories, where it would
    only price the correction it is meant to coexist with. Keep the null constraint as a gauge fix
    and upgrade it to a freeze; it was never single-concept protection. The data term is a different
    project and reaches a quarter of the vocabulary. \emph{Likely but verify.}
  \item \textbf{Off-policy.} The one-step-off-trajectory probe of Section 5: take one adapted DDIM
    step from each cached state, compare the fit loss there against at the cached successor, with the
    bar at a factor of $2$ written into the script first. Coverage against the cross-seed spread of
    cached states is the free precursor and costs no forward passes at all.
\end{enumerate}

\paragraph{The two one-line instrumentation changes this document asks for.}
Log $\norm{\Delta}^2$ beside \texttt{loss\_fit} (\texttt{trainer.py:450--453}), and record the
per-concept factor norms the step-0 additive probe already solves for and throws away
(\texttt{scripts/idea\_probes/additive\_floor\_step0.py}). Neither costs a forward pass, and the
first turns Section 3's bracket into an equality.

\end{document}
